\section{Conclusion}\label{sec:conclusion}

We have identified and addressed a significant deficiency in phase-field Allen-Cahn LBM for simulating droplet dynamics on curved surfaces. The Allen-Cahn sharpening term resists geometric wetting boundary condition corrections, suppressing contact angle accuracy on curved hydrophobic surfaces. We proposed a geometric amplification factor $\alpha_{geo}$ that overcomes this resistance, motivated it from the balance between AC sharpening and geometric wetting, and systematically calibrated it across three parameter dimensions.

The key finding is that geometric amplification is effective only for strongly hydrophobic surfaces ($\theta_{eq} > 120^\circ$), small curvature ratios ($R^* < 3$), and well-resolved interfaces ($D_0/\xi > 3.0$, design guideline). The optimal $\alpha_{geo} = 1.5$ is recommended as a robust default across the tested parameter range for superhydrophobic curved surfaces, bringing the spreading asymmetry ratio to within $0.2\%$ of its grid-converged value $k_\infty \approx 2.66$ (versus $-42.0\%$ without amplification), whose magnitude is consistent with the spreading ratio reported by \citet{liu2015symmetry}. Using this correction, we constructed a comprehensive We $\times$ $R^*$ phase diagram of asymmetric droplet spreading via Allen-Cahn LBM at 828:1 density ratio, revealing three spreading regimes with the asymmetry ratio $k$ peaking at $\mathrm{We} \approx 15$ ($k = 3.30$ at $R^*=1.0$).

The method has been validated across six curvature ratios ($R^* = 0.5$--$3.0$), four contact angles ($\theta_{eq} = 90^\circ$--$162^\circ$), and five grid resolutions ($N = 80$--$180$). At parameters matched to \citet{liu2015symmetry} ($\mathrm{We}=10.6$, $R^*=1.2$, $\theta_{eq}=160^\circ$), the simulated spreading diameter ratio is $k \approx 2.93$ with $\alpha_{geo}=1.5$, qualitatively consistent with the asymmetric spreading reported experimentally; direct quantitative comparison is limited by the different metrics (spreading diameter ratio versus momentum ratio, see Sec.~\ref{sec:experiments}). The method correctly preserves symmetric spreading on concave surfaces ($k = 1.0$), where no asymmetry is physically expected. Grid convergence is achieved at $N \geq 100$, with the asymptotic value $k_\infty \approx 2.66 \pm 0.04$. Full 3D D3Q19 simulations confirm that geometric amplification extends to three-dimensional flows, yielding $k_{max} = 2.47$ at $N=80$ for the primary validation case ($R^*=1.0$, $\mathrm{We}=7.9$, $\theta_{eq}=162^\circ$, $\alpha_{geo}=1.5$), with correct symmetric spreading on flat and convex hemisphere surfaces.

Limitations include: (1) quantitative validation is limited to the cylindrical ridge geometry of Liu et al.; (2) the direct static contact-angle measurement is reported at a single point only ($\theta_{eff} \approx 166^\circ$ with $\alpha=0$ vs $\approx 148^\circ$ with $\alpha=1.5$ at target $162^\circ$, Sec.~\ref{subsec:static}), with $\pm 4^\circ$ accuracy at the diffuse-interface width used here. Notably, $\alpha_{geo}=1.5$ over-corrects the equilibrium angle by $-14^\circ$ in the static setting; we interpret the amplification as a dynamic-impact correction, but a direct contact-angle measurement during impact has not been performed, so the mechanism-level claim rests on the $k_{max}$ evidence; (3) the amplification factor is held constant rather than adapting to local curvature; (4) the recommended default $\alpha_{geo}=1.5$ is calibrated within $\theta_{eq}\leq 162^\circ$, $\xi\approx 4$, moderate $\mathrm{We}$, and over-shoots outside this window ($\theta_{eq}=170^\circ$: $k=4.68$; $\xi=3$: $k=4.30$); and (5) reported $k_{max}$ values are windowed maxima over 2000 steps, and at $\mathrm{We}\geq 10$ the asymmetry is still increasing at the window end. Promising future directions include extending the calibration to three-dimensional hemispherical and concave geometries, developing an adaptive $\alpha_{geo}$ that varies with local curvature and interface thickness, coupling with thermal models for icing applications on curved surfaces, and constructing machine learning surrogates trained on the phase diagram data for real-time prediction of droplet dynamics, and extending the simulation window beyond 2000 steps for the high-We cells, where the asymmetry is still growing at the window end.
